\section{Обеспечение качества и роль QA в проекте}

В проекте QA-функция рассматривалась как непрерывный процесс управления качеством, а не как разовая проверка перед демонстрацией. Роль QA заключалась в том, чтобы поддерживать согласованность между продуктовой идеей, технической реализацией и фактическим пользовательским поведением системы. Практически это выражалось в контроле требований, построении тестовой стратегии, регулярной проверке рисковых сценариев и сопровождении релизного цикла.

\subsection{QA на этапе формирования идеи и требований}

Работа по качеству началась до разработки. На старте проекта контролировалось, чтобы команда фиксировала не только функциональные ожидания, но и проверяемые критерии приемки. Это было важно, поскольку бизнес-игра с несколькими ролями и фазами чувствительна к неоднозначным формулировкам.

Ключевые задачи QA на этом этапе были следующими:

\begin{itemize}
    \item формализация границ MVP и исключение второстепенных функций, не влияющих на основной игровой цикл;
    \item выделение ролей (ведущий, игрок, команда) и ограничений их действий в API;
    \item декомпозиция сквозного сценария на проверяемые шаги: создание игры, подключение участников, запуск проектов, ход дня, ретро-фаза, продолжение;
    \item фиксация рисков, которые могли привести к срыву сессии: неверные переходы фаз, рассинхронизация состояния между командами, ошибки прав доступа;
    \item согласование Definition of Ready для задач, чтобы в разработку попадали только проверяемые истории.
\end{itemize}

Ретроспективно именно эта работа позволила снизить количество переоткрытых задач в середине спринта: команда реже сталкивалась с ситуацией, когда реализация технически завершена, но не соответствует исходному пользовательскому ожиданию.

\subsection{Контроль качества в ходе разработки}

На этапе реализации QA-процесс был встроен в ежедневный ритм команды. Контролировалось, чтобы каждая задача проходила минимальный набор этапов: уточнение критериев, реализация, самопроверка разработчика, проверка по чек-листу, регрессионный прогон перед слиянием.

Для этого применялись следующие практики:

\begin{enumerate}
    \item ведение тест-дизайна на уровне сценариев и негативных кейсов одновременно с декомпозицией backend-задач;
    \item проверка изменений на совместимость с уже реализованными фазами игры \texttt{setup}, \texttt{running}, \texttt{retro}, \texttt{finished};
    \item контроль API-контрактов (статусы, структура JSON, тексты ошибок) для предотвращения скрытых интеграционных дефектов frontend/backend;
    \item контроль инвариантов доменной логики: доступность действий только соответствующим ролям, корректное обновление досок команд, неизбыточное создание сущностей;
    \item регрессионная проверка рисковых мест после каждого изменения логики хода дня и управления WIP-лимитами.
\end{enumerate}

\subsection{Тестовая стратегия backend-части}

Для backend-части проекта была реализована стратегия двух уровней автоматизации на \texttt{pytest}:

\begin{itemize}
    \item \textbf{Unit-suite} (\texttt{tests/unit/test\_backend\_api.py}) --- быстрые контрактные проверки отдельных endpoint-ов и инвариантов состояния;
    \item \textbf{WebSocket-suite} (\texttt{tests/unit/test\_websockets.py}) --- проверка протокола сообщений и корректности реакций лобби и игрового сокета;
    \item \textbf{E2E-suite} (\texttt{tests/e2e/test\_game\_flow.py}) --- сквозные сценарии, моделирующие поведение реальных участников игры.
\end{itemize}

Unit-уровень покрывает прежде всего корректность реакций API на стандартные и ошибочные действия: создание игры, присоединение участников, правила запуска, переходы по дням, ограничения прав, валидацию параметров. WebSocket-уровень фиксирует контракт сообщений и поведение при ошибках подключения (неизвестный код, отсутствие кода), а также жизненный цикл соединения. E2E-уровень проверяет устойчивость всей цепочки: от подготовки сессии до перехода в ретро и продолжения игрового цикла.

Маркировка тестов реализована через \texttt{pytest.ini} с маркерами \texttt{unit} и \texttt{e2e}, что позволило быстро выполнять прицельные прогоны. Отдельного CI/CD в проекте не было, поэтому регрессионные прогоны выполнялись локально разработчиками и QA перед слиянием ключевых изменений.

\subsection{Набор контролируемых тест-кейсов}

В рамках QA-процесса контролировались и поддерживались следующие группы тест-кейсов:

\begin{itemize}
    \item \textbf{Жизненный цикл игры}: корректное создание, запуск проекта, старт игры, переходы между днями и фазами;
    \item \textbf{Ролевая модель}: только ведущий может выполнять управляющие действия (запуск, переход дня, изменение WIP, продолжение после ретро);
    \item \textbf{Валидация и защита от ошибок}: запрет некорректных team/game ID, проверка обязательных полей и диапазонов значений;
    \item \textbf{Согласованность командных досок}: задачи создаются в ожидаемых колонках и не могут перемещаться участником другой команды;
    \item \textbf{Фазовые ограничения}: изменение WIP допускается только в \texttt{retro}, а переход в следующий день --- только после завершения действий команд.
\end{itemize}

Такой набор кейсов позволил закрыть критичные риски регрессии в логике игры и обеспечить предсказуемое поведение API при расширении функциональности.

\subsection{Примеры автоматизированных сценариев}

Ниже приведены примеры автотестов, отражающие ключевые риски:

\begin{itemize}
    \item \textbf{Контракт создания игры}: проверка кода игры, идентификатора ведущего, начальной фазы и структуры команд (\texttt{test\_create\_game\_returns\_expected\_structure}).
    \item \textbf{Ролевая модель}: запрет запуска проекта и игры не-ведущим (\texttt{test\_start\_project\_forbidden\_for\_non\_facilitator}, \texttt{test\_start\_game\_forbidden\_for\_non\_facilitator}).
    \item \textbf{Фазовые ограничения}: запрет перемещения карточек до старта и в ретро-фазе (\texttt{test\_drag\_before\_game\_started}, \texttt{test\_drag\_in\_retro\_phase}).
    \item \textbf{Сквозной цикл}: переходы по дням, ретро-фаза, изменение WIP и продолжение игры (\texttt{test\_full\_game\_cycle\_to\_retro\_and\_continue}).
    \item \textbf{WebSocket-контракты}: корректный ответ на ping/pong и ошибки неправильного запроса (\texttt{test\_lobby\_ws\_ping\_pong}, \texttt{test\_lobby\_ws\_unknown\_type}).
\end{itemize}

\subsection{Окружения и организация тестирования}

Тестирование выполнялось в двух окружениях: локально на машинах разработчиков и QA, а также на серверной сборке, развёрнутой через \texttt{docker-compose} для проверки сквозного пользовательского потока. Основной фреймворк автоматизации --- \texttt{pytest}; отчёты о покрытии не собирались, приоритетом были сценарии поведения и инварианты доменной логики. Перед демонстрациями выполнялась ручная проверка основных маршрутов, а также быстрый регрессионный прогон unit-тестов.

\subsection{Управление дефектами и результаты}

Для учёта задач и дефектов использовались \texttt{GitHub Issues} и бэклог в \texttt{GitHub Projects}. Отдельной системы баг-трекинга не применялось, но дефекты фиксировались как задачи с пометками о приоритете и серьёзности. Приоритизация проводилась вручную: отдельно отмечались блокирующие/критические дефекты и дефекты высокой срочности, влияющие на игровой сценарий и UX.

Во время полной ручной e2e-сессии всей команды было выявлено около пяти блокирующих/критических дефектов и сопоставимое число дефектов высокого приоритета; дефектов среднего и низкого приоритета обнаруживалось больше в ходе ручных и автоматизированных прогонов. Критические дефекты исправлялись до следующей демонстрации; менее приоритетные оставались в бэклоге как улучшения.

\subsection{Процессы команды под контролем QA}

Помимо технического тестирования контролировались процессные аспекты, влияющие на качество результата:

\begin{itemize}
    \item соблюдение Definition of Done перед переносом задачи в \textit{Done};
    \item прозрачный учёт дефектов и разделение их по приоритету (блокирующие игровой сценарий, критичные по UX, некритичные);
    \item проверка готовности сборки к демо по чек-листу (стабильность основных маршрутов, отсутствие критичных ошибок в сценарии занятия);
    \item участие в ретроспективах с акцентом на причины дефектов процесса: позднее уточнение требований, неполное покрытие пограничных состояний, отсутствие регрессии в рисковых точках;
    \item формирование предложений по улучшениям: раннее уточнение критериев, расширение негативных тестов и фиксация API-контрактов до начала реализации.
\end{itemize}

Итогом стало то, что QA в проекте выполнял роль связующего механизма между идеей продукта, инженерной реализацией и учебным пользовательским сценарием. Это повысило управляемость разработки и снизило вероятность критических сбоев во время проведения бизнес-игры.
